\documentclass[14pt,a4paper]{extarticle}
% =============================================================
% preamble.tex — оформление отчётов по СТП БГУИР 01-2024
% =============================================================

% Шрифты: под pdflatex — Times через tempora + T2A,
% под xelatex/tectonic — системный Times New Roman (запасной вариант STIX Two Text)
\usepackage{iftex}
\ifPDFTeX
  \usepackage[utf8]{inputenc}
  \usepackage[T2A]{fontenc}
  \usepackage{tempora}
\else
  \usepackage{fontspec}
  \IfFontExistsTF{Times New Roman}
    {\setmainfont{Times New Roman}}
    {\setmainfont{STIX Two Text}}
  \IfFontExistsTF{DejaVu Sans Mono}
    {\setmonofont{DejaVu Sans Mono}[Scale=0.82]}
    {\setmonofont{DejaVuSansMono}[Extension=.ttf, BoldFont=*-Bold,
                                  ItalicFont=*-Oblique, Scale=0.82]}
\fi
\usepackage[russian]{babel}

% Поля: слева 30 мм, справа 15 мм, сверху и снизу 20 мм
\usepackage{geometry}
\geometry{left=30mm, right=15mm, top=20mm, bottom=20mm}

\usepackage{setspace}
\singlespacing

\pretolerance=1000
\tolerance=2000
\emergencystretch=3em

\usepackage{indentfirst}
\setlength{\parindent}{1.25cm}

% Номер страницы — снизу справа
\usepackage{fancyhdr}
\pagestyle{fancy}
\fancyhf{}
\fancyfoot[R]{\thepage}
\renewcommand{\headrulewidth}{0pt}
\fancypagestyle{plain}{%
  \fancyhf{}%
  \fancyfoot[R]{\thepage}%
  \renewcommand{\headrulewidth}{0pt}%
}

\usepackage{amsmath,amssymb}
\usepackage{graphicx}
\usepackage{float}
\usepackage{booktabs}
\usepackage{tabularx}
\usepackage{multirow}
\usepackage{longtable}
\usepackage{array}

\usepackage{caption}
\captionsetup[figure]{labelsep=endash, justification=centering,
                      font={normalsize}, name=Рисунок}
\captionsetup[table]{labelsep=endash, justification=raggedright,
                     singlelinecheck=false, font={normalsize}, name=Таблица}

\usepackage{titlesec}
\titleformat{\section}
  {\newpage\normalfont\fontsize{14}{18}\selectfont\bfseries}
  {\thesection}{0.5em}{\MakeUppercase}
\titlespacing*{\section}{1.25cm}{0pt}{14pt}
\titleformat{\subsection}
  {\normalfont\fontsize{14}{18}\selectfont\bfseries}
  {\thesubsection}{0.5em}{}
\titlespacing*{\subsection}{1.25cm}{14pt}{14pt}
\titleformat{\subsubsection}
  {\normalfont\fontsize{14}{18}\selectfont\bfseries}
  {\thesubsubsection}{0.5em}{}
\titlespacing*{\subsubsection}{1.25cm}{14pt}{14pt}

\newcommand{\sectioncenter}[1]{%
  \newpage
  \phantomsection
  \addcontentsline{toc}{section}{#1}
  \begin{center}\textbf{\MakeUppercase{#1}}\end{center}
  \vspace{14pt}
}

\usepackage{tocloft}
\renewcommand{\cftsecfont}{\normalfont}
\renewcommand{\cftsecpagefont}{\normalfont}
\renewcommand{\cftsubsecfont}{\normalfont}
\renewcommand{\cftsubsecpagefont}{\normalfont}
\renewcommand{\cftsecleader}{\cftdotfill{\cftdotsep}}
\renewcommand{\cftdotsep}{1}
% babel переопределяет заголовки после преамбулы, поэтому правим через \addto
\addto\captionsrussian{\renewcommand{\contentsname}{\hfill\textbf{СОДЕРЖАНИЕ}\hfill}}
\renewcommand{\cftaftertoctitle}{\hfill}

\usepackage{enumitem}
\setlist[itemize]{noitemsep, topsep=0pt, leftmargin=1.25cm, label=--}
\setlist[enumerate]{noitemsep, topsep=0pt, leftmargin=1.25cm}

% Листинги кода и вывода программы
\usepackage{listings}
\usepackage{fancyvrb}
\lstset{
  basicstyle=\ttfamily\footnotesize,
  frame=single,
  breaklines=true,
  showstringspaces=false,
  inputencoding=utf8,
  extendedchars=true,
  columns=fullflexible,
  keepspaces=true,
  captionpos=t
}
\ifPDFTeX
\lstset{
  literate={а}{{\cyra}}1 {б}{{\cyrb}}1 {в}{{\cyrv}}1 {г}{{\cyrg}}1
           {д}{{\cyrd}}1 {е}{{\cyre}}1 {ж}{{\cyrzh}}1 {з}{{\cyrz}}1
           {и}{{\cyri}}1 {й}{{\cyrishrt}}1 {к}{{\cyrk}}1 {л}{{\cyrl}}1
           {м}{{\cyrm}}1 {н}{{\cyrn}}1 {о}{{\cyro}}1 {п}{{\cyrp}}1
           {р}{{\cyrr}}1 {с}{{\cyrs}}1 {т}{{\cyrt}}1 {у}{{\cyru}}1
           {ф}{{\cyrf}}1 {х}{{\cyrh}}1 {ц}{{\cyrc}}1 {ч}{{\cyrch}}1
           {ш}{{\cyrsh}}1 {щ}{{\cyrshch}}1 {ъ}{{\cyrhrdsn}}1
           {ы}{{\cyrery}}1 {ь}{{\cyrsftsn}}1 {э}{{\cyrerev}}1
           {ю}{{\cyryu}}1 {я}{{\cyrya}}1 {ё}{{\cyryo}}1
           {А}{{\CYRA}}1 {Б}{{\CYRB}}1 {В}{{\CYRV}}1 {Г}{{\CYRG}}1
           {Д}{{\CYRD}}1 {Е}{{\CYRE}}1 {Ж}{{\CYRZH}}1 {З}{{\CYRZ}}1
           {И}{{\CYRI}}1 {Й}{{\CYRISHRT}}1 {К}{{\CYRK}}1 {Л}{{\CYRL}}1
           {М}{{\CYRM}}1 {Н}{{\CYRN}}1 {О}{{\CYRO}}1 {П}{{\CYRP}}1
           {Р}{{\CYRR}}1 {С}{{\CYRS}}1 {Т}{{\CYRT}}1 {У}{{\CYRU}}1
           {Ф}{{\CYRF}}1 {Х}{{\CYRH}}1 {Ц}{{\CYRC}}1 {Ч}{{\CYRCH}}1
           {Ш}{{\CYRSH}}1 {Щ}{{\CYRSHCH}}1 {Ъ}{{\CYRHRDSN}}1
           {Ы}{{\CYRERY}}1 {Ь}{{\CYRSFTSN}}1 {Э}{{\CYREREV}}1
           {Ю}{{\CYRYU}}1 {Я}{{\CYRYA}}1 {Ё}{{\CYRYO}}1
           {—}{{---}}1 {–}{{--}}1 {…}{{\ldots}}1
}
\fi

% Подписи листингов — в том же стиле, что таблицы: «Листинг X.Y — Название»
\captionsetup[lstlisting]{labelsep=endash, justification=raggedright,
                          singlelinecheck=false, font={normalsize}}
\renewcommand{\lstlistingname}{Листинг}
\AtBeginDocument{\numberwithin{lstlisting}{section}}

\lstdefinestyle{python}{language=Python, keywordstyle=\bfseries}
\lstdefinestyle{shell}{language=bash, keywordstyle=\mdseries}
\lstdefinestyle{plain}{language={}, basicstyle=\ttfamily\scriptsize}

% Абзац после листинга должен сохранять отступ 1,25 см (СТП 2.1.1)
\makeatletter
\AtBeginDocument{\let\@endpetrue\@endpefalse}
\makeatother

% Подпись к блоку verbatim: \captionof вне плавающего окружения
% глобально обнуляет \parindent, поэтому отступ восстанавливается вручную
\newcommand{\listingcaption}[2]{%
  \captionof{lstlisting}{#1}\label{#2}%
  \setlength{\parindent}{1.25cm}%
}

\usepackage{hyperref}
\hypersetup{colorlinks=true, linkcolor=black, urlcolor=blue, citecolor=black}

\numberwithin{figure}{section}
\numberwithin{table}{section}
\numberwithin{equation}{section}

% Титульный лист отчёта по лабораторной работе
\newcommand{\makelabtitle}[2]{%
  \begin{titlepage}
  \thispagestyle{empty}
  \begin{center}

  МИНИСТЕРСТВО ОБРАЗОВАНИЯ РЕСПУБЛИКИ БЕЛАРУСЬ

  \vspace{0.5em}
  Учреждение образования

  \vspace{0.5em}
  \textbf{<<Белорусский государственный университет информатики и радиоэлектроники>>}

  \vspace{2em}
  Факультет компьютерных систем и сетей

  \vspace{0.5em}
  Кафедра программного обеспечения информационных технологий
  \end{center}

  \vspace{3em}
  \begin{center}
  \textbf{ОТЧЕТ}

  \vspace{0.5em}
  по лабораторной работе №#1

  \vspace{0.5em}
  по дисциплине

  \vspace{0.5em}
  <<Информационные и технические средства защиты авторских прав>>

  \vspace{0.5em}
  <<#2>>
  \end{center}

  \vspace{5em}
  \begin{flushright}
  \begin{minipage}{0.5\textwidth}
  \textbf{Выполнил:}\\
  Лебедевич Артем Владимирович\\
  магистрант группы 556301

  \vspace{0.6em}
  \textbf{Преподаватель:}\\
  Ярмолик Вячеслав Николаевич\\
  доктор технических наук\\
  профессор кафедры ПОИТ
  \end{minipage}
  \end{flushright}

  \vfill
  \begin{center}Минск, 2026~г.\end{center}
  \end{titlepage}
}

% Титульный лист реферата
\newcommand{\makerefattitle}[1]{%
  \begin{titlepage}
  \thispagestyle{empty}
  \begin{center}

  МИНИСТЕРСТВО ОБРАЗОВАНИЯ РЕСПУБЛИКИ БЕЛАРУСЬ

  \vspace{0.5em}
  Учреждение образования

  \vspace{0.5em}
  \textbf{<<Белорусский государственный университет информатики и радиоэлектроники>>}

  \vspace{2em}
  Факультет компьютерных систем и сетей

  \vspace{0.5em}
  Кафедра программного обеспечения информационных технологий
  \end{center}

  \vspace{3em}
  \begin{center}
  \textbf{РЕФЕРАТ}

  \vspace{0.5em}
  по дисциплине

  \vspace{0.5em}
  <<Информационные и технические средства защиты авторских прав>>

  \vspace{1.5em}
  на тему

  \vspace{0.5em}
  \textbf{<<#1>>}
  \end{center}

  \vspace{5em}
  \begin{flushright}
  \begin{minipage}{0.5\textwidth}
  \textbf{Выполнил:}\\
  Лебедевич Артем Владимирович\\
  магистрант группы 556301

  \vspace{0.6em}
  \textbf{Преподаватель:}\\
  Ярмолик Вячеслав Николаевич\\
  доктор технических наук\\
  профессор кафедры ПОИТ
  \end{minipage}
  \end{flushright}

  \vfill
  \begin{center}Минск, 2026~г.\end{center}
  \end{titlepage}
}


\begin{document}

\makelabtitle{4}{Эмулятор банкомата}

\tableofcontents

\section{Цель работы}

Разработать эмулятор банкомата, взаимодействующий с виртуальным банком из предыдущих
работ: реализовать клиентский интерфейс, повторяющий работу с реальным устройством, и
протокол общения банкомата с банком на основе транзакций. Рабочим счётом карты служит
кредитный счёт, открытый в третьей лабораторной работе.

\section{Задание}

Банкомат --- автоматизированное устройство, которое позволяет держателю счёта удалённо
выполнять операции:
\begin{itemize}
\item аутентификация по номеру карты и PIN-коду; после трёх неверных PIN-кодов подряд
работа с банкоматом принудительно завершается;
\item снятие денег с кредитного счёта;
\item просмотр остатка кредитного счёта;
\item просмотр остатка депозитного счёта (необязательная операция);
\item платёж за услуги оператора мобильной связи (необязательная операция);
\item возврат карты.
\end{itemize}

Модель общения с банком построена на транзакциях. В ходе диалога с пользователем банкомат
накапливает вводимую информацию во \textit{внутреннем списке}: номер карты, PIN-код, вид
операции, её параметры. Когда собраны все данные, необходимые для операции, содержимое
списка целиком отправляется серверу банка, а список очищается. Для следующей операции
данные карты вводятся заново; номер карты эмулятор запоминает, поскольку он имитирует саму
карту, оставшуюся в устройстве.

Операции со счётом сопровождаются чеком: при платеже чек печатается автоматически, в
остальных случаях --- по желанию клиента. Номер карты и PIN-код вводятся с клавиатуры; сумма
снятия должна быть неотрицательным целым числом; номер телефона при платеже состоит из
десяти цифр и перед отправкой подтверждается клиентом. В работе реализованы все операции,
включая необязательные. Дополнительно банкомат работает на трёх языках --- русском,
английском и белорусском, а для знакомства с ним банк выпускает демонстрационные карты.

\section{Проектирование}

\subsection{Архитектура}

Задание выделяет в программе две системы --- клиентский интерфейс и протокол общения
банка и банкомата. Им соответствуют два уровня решения (рисунок~\ref{fig:arch}).

Банкомат реализован отдельным микросервисом \texttt{atm-service}. Он не имеет базы данных
и ничего не знает о счетах и договорах: его задача --- вести диалог с клиентом и собирать
транзакцию. Web-интерфейс выполнен как <<тонкий>> терминал: браузер только рисует экран,
который описал сервис, и сообщает о нажатиях. Вся логика сеанса находится на сервере, где её
можно тестировать без браузера.

Банковскую сторону протокола реализует сервис кредитов --- владелец карты и кредитного
счёта. Он проверяет карту и PIN-код, проводит операции через главную книгу, а остатки
вкладов запрашивает у сервиса депозитов. Банкомату доступна единственная точка входа банка,
и получить доступ к счёту в обход проверки PIN-кода он не может.

\begin{figure}[H]
\centering
\begin{tikzpicture}[lbl/.style={font=\scriptsize, text=black!80}]
\node[box, fill=codeback, minimum width=2.4cm, minimum height=1.2cm] (ui) at (0, 0) {Интерфейс\\[-1pt]банкомата\\[-3pt]{\footnotesize Vue 3}};
\node[svc, fill=codekw!18, minimum width=3.0cm] (atm) at (4.4, 0) {atm-service\\[-1pt]{\footnotesize автомат сеанса,}\\[-3pt]{\footnotesize внутренний список}};
\node[svc, minimum width=3.0cm] (crd) at (9.2, 0) {credit-service\\[-1pt]{\footnotesize карты, PIN-коды,}\\[-3pt]{\footnotesize процессинг}};
\node[svc, minimum width=2.9cm] (acc) at (13.5, 1.3) {account-service\\[-1pt]{\footnotesize проводки, остатки}};
\node[svc, minimum width=2.9cm] (dep) at (13.5, -1.3) {deposit-service\\[-1pt]{\footnotesize вклады клиента}};
\draw[flow, <->] (ui) -- node[above, lbl] {экран} node[below, lbl] {события} (atm);
\draw[flow, <->] (atm) -- node[above, lbl] {транзакция} node[below, lbl] {ответ} (crd);
\draw[flow] (crd.east) -- ++(0.4, 0) |- (acc.west);
\draw[flow] (crd.east) -- ++(0.4, 0) |- (dep.west);
\draw[black!45, dashed] (6.8, 2.1) -- (6.8, -2.1);
\node[note] at (2.6, 2.0) {банкомат};
\node[note] at (10.6, 2.0) {банк};
\end{tikzpicture}
\caption{Банкомат и банк: уровни решения}
\label{fig:arch}
\end{figure}

\subsection{Протокол общения банка и банкомата}

Транзакция --- упорядоченный список пар <<имя --- значение>>, то есть содержимое
внутреннего списка банкомата. Она передаётся одним запросом
\texttt{POST /api/atm/transactions}; элементы списка перечислены в
таблице~\ref{tab:items}. Банк отвечает статусом \texttt{OK} или \texttt{ERROR}, кодом, текстом
для экрана и данными для чека (таблица~\ref{tab:codes}). Пример обмена приведён в
листинге~\ref{lst:exchange}.

\begin{table}[H]
\caption{Элементы транзакции}
\label{tab:items}
\centering
\small
\begin{tabularx}{\textwidth}{llX}
\toprule
Имя & Значение & Когда добавляется в список \\
\midrule
\texttt{CARD} & номер карты, 16 цифр & после ввода номера; перед новой операцией --- автоматически \\
\texttt{PIN} & PIN-код, 4 цифры & после ввода PIN-кода \\
\texttt{OPERATION} & \texttt{AUTHORIZE}, \texttt{WITHDRAW}, \texttt{BALANCE}, \texttt{DEPOSIT\_BALANCE}, \texttt{PAYMENT} & при выборе пункта меню \\
\texttt{AMOUNT} & сумма & после ввода суммы снятия или платежа \\
\texttt{OPERATOR} & код оператора связи & после выбора оператора \\
\texttt{PHONE} & номер телефона, 10 цифр & после ввода номера \\
\bottomrule
\end{tabularx}
\end{table}

\begin{table}[H]
\caption{Коды ответа банка}
\label{tab:codes}
\centering
\small
\begin{tabularx}{\textwidth}{lX}
\toprule
Код & Значение \\
\midrule
\texttt{OK} & операция выполнена; в данных --- сумма, остаток, реквизиты для чека \\
\texttt{WRONG\_PIN} & неверный PIN-код; в данных --- число оставшихся попыток \\
\texttt{CARD\_BLOCKED} & карта заблокирована после трёх неверных PIN-кодов подряд \\
\texttt{CARD\_NOT\_FOUND} & карта не обслуживается банком \\
\texttt{INSUFFICIENT\_FUNDS} & на счёте недостаточно средств \\
\texttt{INVALID\_AMOUNT}, \texttt{INVALID\_PHONE} & неверная сумма или номер телефона \\
\texttt{UNKNOWN\_OPERATOR}, \texttt{UNKNOWN\_OPERATION} & оператора связи или операции нет в справочнике \\
\texttt{SERVICE\_UNAVAILABLE} & смежный сервис банка временно недоступен \\
\bottomrule
\end{tabularx}
\end{table}

\begin{lstlisting}[style=json, caption={Транзакция снятия наличных и ответ банка}, label={lst:exchange}]
POST /api/atm/transactions
{"items": [
  {"name": "CARD",      "value": "9112380000000010"},
  {"name": "PIN",       "value": "****"},
  {"name": "OPERATION", "value": "WITHDRAW"},
  {"name": "AMOUNT",    "value": "450"}
]}

200 OK
{"status": "OK", "code": "OK", "message": "Заберите деньги",
 "data": {"amount": 450, "balance": 4550.00, "currency": "BYN",
          "bankDate": "2026-10-01", "holder": "Лебедевич Артем Владимирович"}}
\end{lstlisting}

Банк не хранит состояние диалога: каждая транзакция содержит номер карты и PIN-код и
проверяется заново. Поэтому банковская сторона остаётся сервисом без состояния и может
масштабироваться горизонтально, а состояние сеанса живёт только в банкомате.

\subsection{Конечный автомат сеанса}

Работа с банкоматом описана конечным автоматом из одиннадцати состояний
(рисунок~\ref{fig:fsm}). События автомата --- ввод значения с клавиатуры (клавиша <<Ввод>>),
выбор пункта меню боковой кнопкой и клавиша <<Отмена>>.

\begin{figure}[H]
\centering
\begin{tikzpicture}[
  st/.style={draw=black!70, rounded corners=4pt, fill=codekw!6, minimum width=3.0cm, minimum height=0.95cm,
             align=center, font=\footnotesize, inner sep=3pt},
  lbl/.style={font=\scriptsize, text=black!85, align=center}]
\node[st] (ic) at (0, 0) {\texttt{INSERT\_CARD}\\ожидание карты};
\node[st] (pin) at (5.7, 0) {\texttt{PIN}\\ввод PIN-кода};
\node[st] (menu) at (11.4, 0) {\texttt{MENU}\\главное меню};
\node[st] (result) at (0, -2.7) {\texttt{RESULT}\\состояние счёта};
\node[st] (receipt) at (5.7, -2.7) {\texttt{RECEIPT\_PROMPT}\\печатать чек?};
\node[st] (amount) at (11.4, -2.7) {\texttt{AMOUNT}\\сумма снятия};
\node[st] (message) at (0, -5.4) {\texttt{MESSAGE}\\итог операции};
\node[st] (confirm) at (5.7, -5.4) {\texttt{CONFIRM}\\подтверждение};
\node[st] (operator) at (11.4, -5.4) {\texttt{OPERATOR}\\выбор оператора};
\node[st] (payamount) at (5.7, -8.1) {\texttt{PAY\_AMOUNT}\\сумма платежа};
\node[st] (phone) at (11.4, -8.1) {\texttt{PHONE}\\номер телефона};

\draw[flow] ([yshift=2mm]ic.east) -- node[above, lbl] {номер карты} ([yshift=2mm]pin.west);
\draw[flow] ([yshift=-2mm]pin.west) -- node[below, lbl] {3 ошибки} ([yshift=-2mm]ic.east);
\draw[flow] ([yshift=2mm]pin.east) -- node[above, lbl] {PIN-код принят} ([yshift=2mm]menu.west);
\draw[flow] ([yshift=-2mm]menu.west) -- node[below, lbl] {новая операция} ([yshift=-2mm]pin.east);
\draw[flow] (pin) edge[loop below, looseness=6] node[lbl, below] {неверный PIN-код} (pin);
\draw[flow] (menu) -- node[right, lbl] {снять\\наличные} (amount);
\draw[flow] (amount) -- node[above, lbl] {деньги выданы} (receipt);
\draw[flow, <->] (menu.south west) -- node[lbl, sloped, above] {остаток счёта} node[lbl, sloped, below] {после снятия} (receipt.north east);
\draw[flow] (receipt) -- node[above, lbl] {после остатка} (result);
\draw[flow] (result) -- node[right, lbl] {забрать\\карту} (ic);
\draw[flow] (menu.east) -- ++(0.75, 0) |- (operator.east);
\node[lbl, anchor=east] at ([xshift=6mm, yshift=-13.5mm]amount.east) {оплата связи};
\draw[flow] (operator) -- node[right, lbl] {оператор} (phone);
\draw[flow] (phone) -- node[above, lbl] {номер} (payamount);
\draw[flow] (payamount) -- node[right, lbl] {сумма} (confirm);
\draw[flow] (confirm) -- node[above, lbl] {ввести заново} (operator);
\draw[flow] (confirm) -- node[above, lbl] {подтвердить} (message);
\draw[flow] (message.west) -- ++(-0.75, 0) |- ([xshift=6mm, yshift=13mm]menu.north) -- ([xshift=6mm]menu.north);
\draw[black!75, line width=0.5pt] (result.west) -- ++(-0.75, 0);
\node[lbl] at (5.7, 2.0) {продолжить работу};
\draw[flow] ([xshift=-6mm]menu.north) -- ++(0, 0.55) -| (ic.north);
\node[lbl] at (5.7, 1.25) {забрать карту};
\end{tikzpicture}
\caption{Конечный автомат сеанса работы с банкоматом (основные переходы)}
\label{fig:fsm}
\end{figure}

На рисунке не показаны переходы по отказу банка: из состояний \texttt{AMOUNT} и
\texttt{CONFIRM}, а также при запросе остатка отказ (например, нехватка средств) ведёт в
состояние \texttt{MESSAGE}, откуда клиент возвращается в меню. Клавиша <<Отмена>> внутри
операции возвращает в меню, а в меню и при вводе PIN-кода --- возвращает карту.

\subsection{Внутренний список}

Жизненный цикл внутреннего списка на примере двух операций подряд показан в
таблице~\ref{tab:buffer}. После авторизации номер карты и PIN-код остаются в списке и
используются первой операцией. Отправка транзакции очищает список, поэтому перед второй
операцией банкомат сам добавляет номер карты, запрашивает PIN-код и повторно авторизует
карту --- неверный PIN-код обнаруживается сразу, а не после ввода суммы.

\begin{table}[H]
\caption{Внутренний список при снятии наличных и последующем запросе остатка}
\label{tab:buffer}
\centering
\small
\begin{tabularx}{\textwidth}{lXl}
\toprule
Действие клиента & Содержимое списка & Обращение к банку \\
\midrule
Ввод номера карты & \texttt{CARD} & --- \\
Ввод PIN-кода & \texttt{CARD, PIN} & \texttt{AUTHORIZE} \\
Выбор <<Снять наличные>> & \texttt{CARD, PIN, OPERATION} & --- \\
Ввод суммы & \texttt{CARD, PIN, OPERATION, AMOUNT} & \texttt{WITHDRAW} \\
Ответ банка получен & пуст & --- \\
Выбор <<Остаток счёта>> & \texttt{CARD} (подставлен банкоматом) & --- \\
Ввод PIN-кода & \texttt{CARD, PIN} & \texttt{AUTHORIZE} \\
 & \texttt{CARD, PIN, OPERATION} & \texttt{BALANCE} \\
Ответ банка получен & пуст & --- \\
\bottomrule
\end{tabularx}
\end{table}

\subsection{Операции банка}

Снятие наличных в банкомате --- это операции 2 и 3 схемы <<Кредитная программа>>:
<<Перевод кредита в кассу>> (кредит текущего счёта, дебет кассы) и <<Получение кредита через
кассу>> (кредит кассы). Перед проводкой сервис кредитов сравнивает сумму с остатком текущего
счёта; тот же запрет отрицательного остатка действует и в главной книге.

Платёж за услуги связи списывает сумму с текущего счёта клиента и зачисляет её на
расчётный счёт оператора --- пассивный счёт 3012 <<Текущие (расчётные) счета коммерческих
организаций>>, который открывается в главной книге при первом платеже этому оператору.

Запрос остатка депозитного счёта сервис кредитов переадресует сервису депозитов и
возвращает действующие вклады держателя карты.

\subsection{Защита PIN-кода}

PIN-код генерируется при выпуске карты случайным образом и показывается сотруднику
банка один раз (<<ПИН-конверт>>). В базе хранится только его хеш BCrypt, восстановить
PIN-код по которому нельзя; при утрате он перевыпускается. Счётчик неверных попыток
ведёт банк, а не банкомат: после третьей ошибки подряд карта блокируется, и начать подбор
заново в новом сеансе или на другом банкомате невозможно. Блокировка снимается
перевыпуском PIN-кода. На экране, во внутреннем списке и в журнале обмена PIN-код
отображается звёздочками.

Номер карты содержит контрольную цифру по алгоритму Луна. Банкомат проверяет её сам и
не обращается в банк, если номер набран с ошибкой.

\subsection{Демонстрационные карты}

PIN-код обычной карты известен только тому, кто видел <<ПИН-конверт>>, поэтому сразу после
запуска системы работать с банкоматом было бы нечем. Для знакомства с эмулятором сервис
кредитов после старта сам заключает три кредитных договора и назначает их картам PIN-коды из
настройки \texttt{bank.demo} (таблица~\ref{tab:demo}). Номера и PIN-коды этих карт показаны на
странице банкомата, карту можно <<вставить>> одной кнопкой.

\begin{table}[H]
\caption{Демонстрационные карты}
\label{tab:demo}
\centering
\small
\begin{tabular}{llllr}
\toprule
Номер карты & PIN-код & Держатель & Договор & Кредит, BYN \\
\midrule
9112 3800 0000 0010 & 1111 & Лебедевич А.~В. & К-900001 & 5\,000 \\
9112 3800 0000 0028 & 2222 & Савицкая М.~О. & К-900002 & 3\,000 \\
9112 3800 0000 0036 & 3333 & Жуков Д.~А. & К-900003 & 10\,000 \\
\bottomrule
\end{tabular}
\end{table}

Демо-договоры заключаются обычным путём --- тем же методом, что и договоры из формы, с теми
же счетами и проводками, а PIN-код и здесь хранится только в виде хеша. Сервисам клиентов и
счетов нужно время на запуск, поэтому выдача выполняется в фоновом потоке с повторами, пока
смежные сервисы не ответят; повторный запуск ничего не дублирует. Номера демо-договоров
начинаются с К-900001 и не занимают номера обычных договоров. Демо-режим отключается
параметром \texttt{bank.demo.enabled=false}: в рабочей системе известных заранее PIN-кодов
быть не должно.

\subsection{Локализация}

Банкомат --- <<тонкий>> терминал: тексты экрана, пункты меню и чеки формирует сервис, а не
браузер. Поэтому переводятся они на сервере. Web-интерфейс передаёт выбранный язык в
заголовке \texttt{Accept-Language}, автомат сеанса строит экран на этом языке, а при
обращении к банку заголовок передаётся дальше, и сообщение банка (например, <<Недостаточно
средств>>) приходит на том же языке. В словаре самого интерфейса остаются только подписи
клавиш и служебной панели.

\section{Реализация}

\subsection{Структура}

Состав модуля банкомата и банковской части протокола приведён в листинге~\ref{lst:tree}.

\listingcaption{Модуль \texttt{atm-service} и пакет протокола в \texttt{credit-service}}{lst:tree}
\begin{Verbatim}[fontsize=\footnotesize, frame=single, rulecolor=\color{codeframe}]
atm-service/
  bank/        BankGateway — REST-клиент банка; Transaction, Item, BankReply
  session/     AtmSession  — конечный автомат сеанса и внутренний список
               AtmState, Screen, Receipts, AtmService (реестр сеансов)
  web/         AtmController — REST клиентского интерфейса
  resources/   messages*.properties — тексты экранов и чеков на трёх языках
               static/index.html, app.js, i18n.js, atm.css — интерфейс банкомата
credit-service/
  atm/         AtmProcessor — обработка транзакций; AtmRequest, AtmResponse, DepositApi
  domain/      Card, CardNumbers, MobileOperator, MobilePayment
  service/     CardService — выпуск карты, проверка и перевыпуск PIN-кода
               DemoCards, DemoProperties — демонстрационные карты
  web/         AtmController — POST /api/atm/transactions, GET /api/atm/demo-cards
\end{Verbatim}

\subsection{Автомат сеанса и внутренний список}

Сеанс представлен объектом \texttt{AtmSession}. Событие <<Ввод>> обрабатывается в
зависимости от текущего состояния (листинг~\ref{lst:enter}): значение проверяется, добавляется
во внутренний список \texttt{buffer}, и автомат переходит в следующее состояние либо
обращается к банку. Сообщения заданы ключами словаря и выводятся на языке клиента.

\begin{lstlisting}[style=java, caption={Обработка ввода с клавиатуры (\texttt{AtmSession.enter})}, label={lst:enter}]
public synchronized Screen enter(String text) {
    begin();
    String value = text == null ? "" : text.replace(" ", "");
    switch (state) {
        case INSERT_CARD -> insertCard(value);
        case PIN -> {
            if (!value.matches("\\d{4}")) {
                warn("notice.pinLength");                  // «PIN-код состоит из 4 цифр»
            } else {
                buffer.add(new Item(Item.PIN, value));
                authorize();
            }
        }
        case AMOUNT -> {
            if (acceptAmount(value)) {
                execute();
            }
        }
        case PHONE -> {
            if (!value.matches("\\d{10}")) {
                warn("notice.phoneLength");                // «Номер телефона — 10 цифр…»
            } else {
                buffer.add(new Item(Item.PHONE, value));
                state = AtmState.PAY_AMOUNT;
            }
        }
        case PAY_AMOUNT -> {
            if (acceptAmount(value)) {
                state = AtmState.CONFIRM;
            }
        }
        default -> { }
    }
    return screen();
}
\end{lstlisting}

Отправка транзакции и очистка списка сосредоточены в одном методе
(листинг~\ref{lst:execute}). По ответу банка автомат выбирает следующее состояние: при
платеже чек печатается сразу, при снятии и запросе остатка клиенту задаётся вопрос о чеке,
при отказе показывается сообщение банка.

\begin{lstlisting}[style=java, caption={Отправка транзакции банку (\texttt{AtmSession.execute})}, label={lst:execute}]
/** Отправка накопленной транзакции: содержимое списка уходит банку целиком, список очищается. */
private void execute() {
    reply = send(new ArrayList<>(buffer));
    buffer.clear();
    if (reply == null) {
        state = AtmState.MESSAGE;
        warn("notice.noLinkOperation");                    // «Нет связи с банком. Операция не выполнена»
    } else if (!reply.ok()) {
        if ("WRONG_PIN".equals(reply.code()) || "CARD_BLOCKED".equals(reply.code())) {
            eject(Messages.get("notice.takeCard", reply.message()), false);
        } else {
            state = AtmState.MESSAGE;
            notice = reply.message();
        }
    } else if ("PAYMENT".equals(operation)) {
        receipt = Receipts.print(operation, reply, card, terminalId, clock);     // чек платежа печатается всегда
        state = AtmState.MESSAGE;
    } else {
        if ("WITHDRAW".equals(operation)) {
            cash = Receipts.money(reply, "amount");
        }
        state = AtmState.RECEIPT_PROMPT;
    }
}
\end{lstlisting}

Правило <<данные карты вводятся заново>> реализовано при выборе операции
(листинг~\ref{lst:choose}): если в списке нет номера карты и PIN-кода, банкомат подставляет
запомненный номер и переходит к вводу PIN-кода, запомнив выбранную операцию.

\begin{lstlisting}[style=java, caption={Повторный ввод данных карты (\texttt{AtmSession.chooseOperation})}, label={lst:choose}]
private void chooseOperation(String option) {
    if (hasCredentials()) {
        startOperation(option);
    } else {
        buffer.clear();
        buffer.add(new Item(Item.CARD, card));
        pendingOperation = option;
        state = AtmState.PIN;
    }
}
\end{lstlisting}

\subsection{Банковская сторона протокола}

Транзакцию принимает класс \texttt{AtmProcessor} сервиса кредитов
(листинг~\ref{lst:processor}). Проверка карты и PIN-кода выполняется для каждой транзакции
до разбора операции. Отказы возвращаются не исключениями, а обычным ответом протокола: это
важно для счётчика неверных попыток --- он изменяется в той же транзакции базы данных и
должен быть сохранён.

\begin{lstlisting}[style=java, caption={Проверка карты и PIN-кода, выбор операции (\texttt{AtmProcessor.handle})}, label={lst:processor}]
Card card = cards.findByNumber(tx.get("CARD")).orElse(null);
if (card == null) {
    return AtmResponse.error("CARD_NOT_FOUND", Messages.get("atm.cardNotFound"));
}
if (card.isBlocked()) {
    return AtmResponse.error("CARD_BLOCKED", Messages.get("atm.cardBlocked"));
}
if (!cardService.matches(card, tx.get("PIN"))) {
    int left = card.registerFailedAttempt();
    return left == 0
            ? AtmResponse.error("CARD_BLOCKED", Messages.get("atm.blockedNow"))
            : AtmResponse.error("WRONG_PIN", Messages.get("atm.wrongPin", left), Map.of("attemptsLeft", left));
}
card.setFailedAttempts(0);

try {
    return switch (operation) {
        case "AUTHORIZE" -> AtmResponse.ok(Messages.get("atm.authorized"), Map.of("holder", card.getHolder()));
        case "BALANCE" -> balance(card);
        case "DEPOSIT_BALANCE" -> depositBalance(card);
        case "WITHDRAW" -> withdraw(card, tx.get("AMOUNT"));
        case "PAYMENT" -> payment(card, tx.get("OPERATOR"), tx.get("PHONE"), tx.get("AMOUNT"));
        default -> AtmResponse.error("UNKNOWN_OPERATION", Messages.get("atm.unknownOperation"));
    };
} catch (BankException e) {
    return AtmResponse.error(e.getCode(), e.getMessage());
} ...
\end{lstlisting}

Банк не доверяет проверкам банкомата и повторяет их: сумма снятия должна быть целым
положительным числом, номер телефона --- десятью цифрами с кодом оператора 25, 29, 33 или 44,
оператор --- записью справочника.

\subsection{Клиентский интерфейс}

Интерфейс повторяет лицевую панель банкомата: экран, по четыре боковые кнопки с каждой
стороны, цифровая клавиатура с клавишами <<Отмена>>, <<Сброс>> и <<Ввод>>, приёмник карты,
лотки выдачи денег и чека. Сервис возвращает описание экрана --- заголовок, строки текста,
вид ожидаемого ввода, пункты меню, выданные деньги, строки чека; интерфейс отображает его
и отправляет одно из трёх событий (таблица~\ref{tab:rest}). Пункты меню привязаны к боковым
кнопкам. Пока идёт обращение к банку, на экране показывается надпись <<Авторизация на
сервере...>>. Цифры можно набирать и с клавиатуры компьютера.

Справа от банкомата для наглядности выводятся демонстрационные карты, текущее состояние
автомата, содержимое внутреннего списка и последняя транзакция с ответом банка.

\begin{table}[H]
\caption{REST клиентского интерфейса банкомата}
\label{tab:rest}
\centering
\small
\begin{tabularx}{\textwidth}{llX}
\toprule
Метод & Адрес & Назначение \\
\midrule
POST & \texttt{/api/sessions} & начало сеанса; в ответе --- идентификатор и первый экран \\
GET & \texttt{/api/sessions/\{id\}} & текущий экран \\
POST & \texttt{/api/sessions/\{id\}/enter} & клавиша <<Ввод>> с набранным значением \\
POST & \texttt{/api/sessions/\{id\}/select} & боковая кнопка: выбран пункт меню \\
POST & \texttt{/api/sessions/\{id\}/cancel} & клавиша <<Отмена>> \\
GET & \texttt{/api/demo-cards} & демонстрационные карты с PIN-кодами \\
\bottomrule
\end{tabularx}
\end{table}

\section{Тестирование}

Банкомат и банковская сторона протокола покрыты 77 тестами
(таблица~\ref{tab:tests}). Автомат сеанса тестируется как обычный объект с заглушкой банка:
тесты подают события и проверяют экран, содержимое внутреннего списка и транзакцию,
отправленную банку (листинг~\ref{lst:test}). Функциональные тесты проходят сценарии в
браузере через WebDriver. Итог прогона приведён в листинге~\ref{lst:run}.

\begin{table}[H]
\caption{Тесты эмулятора банкомата и протокола}
\label{tab:tests}
\centering
\small
\begin{tabularx}{\textwidth}{lcX}
\toprule
Класс тестов & Тестов & Что проверяется \\
\midrule
\multicolumn{3}{l}{\textit{atm-service}} \\
\texttt{AtmSessionTest} & 28 & переходы автомата, внутренний список, чеки, отказы, языки \\
\texttt{AtmApiIT} & 5 & REST интерфейса, независимость сеансов, демо-карты \\
\texttt{AtmUiIT} & 4 & снятие с чеком, три неверных PIN-кода, платёж, языки --- в браузере \\
\midrule
\multicolumn{3}{l}{\textit{credit-service}} \\
\texttt{AtmProtocolIT} & 25 & все операции протокола, блокировка карты, проводки \\
\texttt{DemoCardsIT} & 5 & выпуск демо-карт, ответы банка на трёх языках \\
\texttt{CardNumbersTest} & 10 & алгоритм Луна, счётчик попыток \\
\midrule
Всего & 77 & \\
\bottomrule
\end{tabularx}
\end{table}

\begin{lstlisting}[style=java, caption={Тест правила повторного ввода PIN-кода (\texttt{AtmSessionTest})}, label={lst:test}]
@Test
@DisplayName("Перед следующей операцией PIN-код вводится заново, номер карты банкомат помнит сам")
void secondOperationRequiresPinAgain() {
    login(atm);
    atm.select("WITHDRAW");
    atm.enter("100");
    atm.select("NO");

    Screen pin = atm.select("BALANCE");
    assertThat(pin.state()).isEqualTo("PIN");
    assertThat(pin.title()).isEqualTo("Введите PIN-код ещё раз");
    assertThat(pin.buffer()).containsExactly(new Item("CARD", CARD));

    Screen prompt = atm.enter(PIN);
    assertThat(prompt.state()).isEqualTo("RECEIPT_PROMPT");
    assertThat(sent().items()).containsExactly(new Item("CARD", CARD), new Item("PIN", PIN), new Item("OPERATION", "BALANCE"));
}
\end{lstlisting}

\listingcaption{Результат выполнения тестов}{lst:run}
\VerbatimInput[fontsize=\fontsize{7}{8.8}\selectfont, frame=single, rulecolor=\color{codeframe}, samepage=true]{listings/tests.txt}

\section{Результаты работы}

Все пять сервисов запущены, банкомат открыт по адресу \texttt{http://localhost:8085}.
Сразу после запуска на странице видны три демонстрационные карты с PIN-кодами
(рисунок~\ref{fig:demo}); для сценария используется первая из них --- карта к кредиту
5\,000 BYN.

\screenshot{img/demo_cards.png}{Банкомат после запуска и демонстрационные карты}{fig:demo}

Сеанс начинается с ввода номера карты (рисунок~\ref{fig:insert}): номер набирается на
клавиатуре либо подставляется кнопкой <<Вставить>> и отображается группами по четыре цифры.
После неверного PIN-кода банкомат сообщает число оставшихся попыток и запрашивает его
повторно (рисунок~\ref{fig:wrong}); в панели справа виден ответ банка с кодом
\texttt{WRONG\_PIN}.

\screenshot{img/insert_card.png}{Ввод номера карты}{fig:insert}

\screenshot{img/wrong_pin.png}{Неверный PIN-код: повторный запрос}{fig:wrong}

После верного PIN-кода открывается главное меню (рисунок~\ref{fig:menu}). Во внутреннем
списке остаются номер карты и скрытый PIN-код, в журнале обмена --- транзакция авторизации.

\screenshot{img/menu.png}{Главное меню}{fig:menu}

Снятие наличных показано на рисунках~\ref{fig:amount}--\ref{fig:receipt}: клиент вводит
сумму, банк подтверждает операцию, деньги появляются в лотке выдачи, и по желанию клиента
печатается чек с суммой и остатком счёта. После отправки транзакции внутренний список пуст.

\screenshot{img/withdraw_amount.png}{Ввод суммы снятия}{fig:amount}

\screenshot{img/withdraw_cash.png}{Деньги выданы, вопрос о печати чека}{fig:cash}

\screenshot{img/withdraw_receipt.png}{Чек о выдаче наличных}{fig:receipt}

Перед следующей операцией банкомат запрашивает PIN-код ещё раз: номер карты уже
подставлен в список (рисунок~\ref{fig:pin-again}). Остаток кредитного счёта выводится на
экран и в чек; на этом шаге клиент может продолжить работу или забрать карту
(рисунок~\ref{fig:balance}).

\screenshot{img/pin_again.png}{Повторный ввод PIN-кода перед новой операцией}{fig:pin-again}

\screenshot{img/balance.png}{Остаток кредитного счёта}{fig:balance}

Оплата мобильной связи: выбор оператора, ввод номера и суммы, экран подтверждения
(рисунок~\ref{fig:confirm}) и автоматически напечатанный чек (рисунок~\ref{fig:payment}).

\screenshot{img/payment_confirm.png}{Подтверждение данных платежа}{fig:confirm}

\screenshot{img/payment_receipt.png}{Платёж принят, чек напечатан автоматически}{fig:payment}

При попытке снять больше, чем доступно на счёте, банк отклоняет транзакцию с кодом
\texttt{INSUFFICIENT\_FUNDS}, и банкомат показывает его сообщение
(рисунок~\ref{fig:insufficient}). Остаток депозитного счёта держателя карты приведён на
рисунке~\ref{fig:deposit}: данные получены сервисом кредитов от сервиса депозитов.

\screenshot{img/insufficient.png}{Отказ банка: недостаточно средств}{fig:insufficient}

\screenshot{img/deposit_balance.png}{Остаток депозитного счёта}{fig:deposit}

После трёх неверных PIN-кодов подряд работа принудительно завершается, а банк блокирует
карту (рисунок~\ref{fig:blocked}): для этого использована третья демонстрационная карта, и
в списке демо-карт она отмечена заблокированной.

\screenshot{img/blocked.png}{Три неверных PIN-кода: карта заблокирована}{fig:blocked}

Банкомат на английском и белорусском языках показан на рисунках~\ref{fig:lang-en}
и~\ref{fig:lang-be}. Переведены экран, пункты меню, клавиши, служебная панель и чек; ответ
банка в журнале обмена также пришёл на выбранном языке.

\screenshot{img/lang_en.png}{Главное меню на английском языке}{fig:lang-en}

\screenshot{img/lang_be.png}{Снятие наличных и чек на белорусском языке}{fig:lang-be}

Результат операций со стороны банка показан на рисунке~\ref{fig:bank}: в карточке
кредитного договора первой карты доступный остаток уменьшился, в проводках видны снятие
через кассу и оплата услуг связи. В отчёте главной книги (рисунок~\ref{fig:report}) появился
расчётный счёт оператора 3012 с суммой платежа.

\screenshot[0.95\textwidth]{img/bank_side.png}{Карточка кредитного договора после операций в банкомате}{fig:bank}

\screenshot{img/report.png}{Отчёт по счетам: расчётный счёт оператора связи}{fig:report}

\sectioncenter{Заключение}

Разработан эмулятор банкомата, работающий с виртуальным банком из предыдущих работ.
Банкомат выполнен отдельным микросервисом: сеанс работы описан конечным автоматом, вводимые
данные накапливаются во внутреннем списке, который целиком отправляется банку при
совершении операции и после этого очищается. Реализованы все операции задания, включая
необязательные: авторизация с ограничением в три попытки, снятие наличных, просмотр остатков
кредитного и депозитного счетов, оплата мобильной связи с подтверждением данных, печать
чеков и возврат карты.

Банковская сторона протокола --- единственная точка входа в сервисе кредитов, которая
при каждой транзакции заново проверяет карту и PIN-код и проводит операции через главную
книгу. Разделение ответственности получилось таким же, как в реальных системах: устройство
ведёт диалог и не имеет доступа к счетам, а банк не хранит состояние диалога и принимает
решения по каждой транзакции независимо. PIN-код хранится только в виде хеша, счётчик
неверных попыток ведётся на стороне банка.

Интерфейс банкомата, экраны и чеки доступны на русском, английском и белорусском языках, а
демонстрационные карты с известными PIN-кодами позволяют опробовать эмулятор сразу после
запуска.

Решение проверено 77 автоматическими тестами: модульными тестами автомата,
интеграционными тестами протокола и функциональными тестами интерфейса через WebDriver.

\sectioncenter{Список использованных источников}

\begin{enumerate}
\item Хопкрофт, Дж. Введение в теорию автоматов, языков и вычислений~/ Дж.~Хопкрофт,
Р.~Мотвани, Дж.~Ульман. --- 2-е изд. --- М.: Вильямс, 2008. --- 528~с.
\item Ньюмен, С. Создание микросервисов~/ С.~Ньюмен. --- 2-е изд. --- СПб.: Питер, 2023. ---
624~с.
\item ISO/IEC 7812-1:2017. Identification cards --- Identification of issuers --- Part 1:
Numbering system. --- Geneva: ISO, 2017. --- 9~p.
\item Provos, N. A Future-Adaptable Password Scheme~/ N.~Provos, D.~Mazi\`eres~// Proceedings
of the USENIX Annual Technical Conference. --- 1999. --- P.~81--91.
\item Spring Framework Reference: REST Clients, HTTP Interface [Электронный ресурс]. ---
Режим доступа: \url{https://docs.spring.io/spring-framework/reference/integration/rest-clients.html}.
--- Дата доступа: 01.10.2026.
\end{enumerate}

\appendixtitle{А}{Исходный код эмулятора банкомата и протокола}

% Файл сгенерирован scripts/gen_listings.py — вручную не править

\subsection*{Эмулятор банкомата (atm-service)}

\lstinputlisting[style=xml, style=appendix, caption={atm-service/pom.xml}]{../../atm-service/pom.xml}
\lstinputlisting[style=yaml, style=appendix, caption={atm-service/src/main/resources/application.yml}]{../../atm-service/src/main/resources/application.yml}
\lstinputlisting[style=java, style=appendix, caption={atm-service/src/main/java/by/bsuir/bank/atm/bank/BankConfig.java}]{../../atm-service/src/main/java/by/bsuir/bank/atm/bank/BankConfig.java}
\lstinputlisting[style=java, style=appendix, caption={atm-service/src/main/java/by/bsuir/bank/atm/bank/BankGateway.java}]{../../atm-service/src/main/java/by/bsuir/bank/atm/bank/BankGateway.java}
\lstinputlisting[style=java, style=appendix, caption={atm-service/src/main/java/by/bsuir/bank/atm/bank/BankReply.java}]{../../atm-service/src/main/java/by/bsuir/bank/atm/bank/BankReply.java}
\lstinputlisting[style=java, style=appendix, caption={atm-service/src/main/java/by/bsuir/bank/atm/bank/Item.java}]{../../atm-service/src/main/java/by/bsuir/bank/atm/bank/Item.java}
\lstinputlisting[style=java, style=appendix, caption={atm-service/src/main/java/by/bsuir/bank/atm/bank/Operator.java}]{../../atm-service/src/main/java/by/bsuir/bank/atm/bank/Operator.java}
\lstinputlisting[style=java, style=appendix, caption={atm-service/src/main/java/by/bsuir/bank/atm/bank/Transaction.java}]{../../atm-service/src/main/java/by/bsuir/bank/atm/bank/Transaction.java}
\lstinputlisting[style=java, style=appendix, caption={atm-service/src/main/java/by/bsuir/bank/atm/session/AtmService.java}]{../../atm-service/src/main/java/by/bsuir/bank/atm/session/AtmService.java}
\lstinputlisting[style=java, style=appendix, caption={atm-service/src/main/java/by/bsuir/bank/atm/session/AtmSession.java}]{../../atm-service/src/main/java/by/bsuir/bank/atm/session/AtmSession.java}
\lstinputlisting[style=java, style=appendix, caption={atm-service/src/main/java/by/bsuir/bank/atm/session/AtmState.java}]{../../atm-service/src/main/java/by/bsuir/bank/atm/session/AtmState.java}
\lstinputlisting[style=java, style=appendix, caption={atm-service/src/main/java/by/bsuir/bank/atm/session/Receipts.java}]{../../atm-service/src/main/java/by/bsuir/bank/atm/session/Receipts.java}
\lstinputlisting[style=java, style=appendix, caption={atm-service/src/main/java/by/bsuir/bank/atm/session/Screen.java}]{../../atm-service/src/main/java/by/bsuir/bank/atm/session/Screen.java}
\lstinputlisting[style=java, style=appendix, caption={atm-service/src/main/java/by/bsuir/bank/atm/AtmProperties.java}]{../../atm-service/src/main/java/by/bsuir/bank/atm/AtmProperties.java}
\lstinputlisting[style=java, style=appendix, caption={atm-service/src/main/java/by/bsuir/bank/atm/AtmServiceApplication.java}]{../../atm-service/src/main/java/by/bsuir/bank/atm/AtmServiceApplication.java}
\lstinputlisting[style=java, style=appendix, caption={atm-service/src/main/java/by/bsuir/bank/atm/web/AtmController.java}]{../../atm-service/src/main/java/by/bsuir/bank/atm/web/AtmController.java}
\lstinputlisting[style=xml, style=appendix, caption={atm-service/src/main/resources/static/index.html}]{../../atm-service/src/main/resources/static/index.html}
\lstinputlisting[style=js, style=appendix, caption={atm-service/src/main/resources/static/app.js}]{../../atm-service/src/main/resources/static/app.js}
\lstinputlisting[style=css, style=appendix, caption={atm-service/src/main/resources/static/atm.css}]{../../atm-service/src/main/resources/static/atm.css}
\lstinputlisting[style=java, style=appendix, caption={atm-service/src/test/java/by/bsuir/bank/atm/AtmApiIT.java}]{../../atm-service/src/test/java/by/bsuir/bank/atm/AtmApiIT.java}
\lstinputlisting[style=java, style=appendix, caption={atm-service/src/test/java/by/bsuir/bank/atm/AtmSessionTest.java}]{../../atm-service/src/test/java/by/bsuir/bank/atm/AtmSessionTest.java}
\lstinputlisting[style=java, style=appendix, caption={atm-service/src/test/java/by/bsuir/bank/atm/AtmUiIT.java}]{../../atm-service/src/test/java/by/bsuir/bank/atm/AtmUiIT.java}

\subsection*{Банковская сторона протокола (credit-service)}

\lstinputlisting[style=java, style=appendix, caption={credit-service/src/main/java/by/bsuir/bank/credit/domain/Card.java}]{../../credit-service/src/main/java/by/bsuir/bank/credit/domain/Card.java}
\lstinputlisting[style=java, style=appendix, caption={credit-service/src/main/java/by/bsuir/bank/credit/domain/CardNumbers.java}]{../../credit-service/src/main/java/by/bsuir/bank/credit/domain/CardNumbers.java}
\lstinputlisting[style=java, style=appendix, caption={credit-service/src/main/java/by/bsuir/bank/credit/domain/MobileOperator.java}]{../../credit-service/src/main/java/by/bsuir/bank/credit/domain/MobileOperator.java}
\lstinputlisting[style=java, style=appendix, caption={credit-service/src/main/java/by/bsuir/bank/credit/domain/MobilePayment.java}]{../../credit-service/src/main/java/by/bsuir/bank/credit/domain/MobilePayment.java}
\lstinputlisting[style=java, style=appendix, caption={credit-service/src/main/java/by/bsuir/bank/credit/dto/PinEnvelope.java}]{../../credit-service/src/main/java/by/bsuir/bank/credit/dto/PinEnvelope.java}
\lstinputlisting[style=java, style=appendix, caption={credit-service/src/main/java/by/bsuir/bank/credit/service/CardService.java}]{../../credit-service/src/main/java/by/bsuir/bank/credit/service/CardService.java}
\lstinputlisting[style=java, style=appendix, caption={credit-service/src/main/java/by/bsuir/bank/credit/atm/AtmConfig.java}]{../../credit-service/src/main/java/by/bsuir/bank/credit/atm/AtmConfig.java}
\lstinputlisting[style=java, style=appendix, caption={credit-service/src/main/java/by/bsuir/bank/credit/atm/AtmProcessor.java}]{../../credit-service/src/main/java/by/bsuir/bank/credit/atm/AtmProcessor.java}
\lstinputlisting[style=java, style=appendix, caption={credit-service/src/main/java/by/bsuir/bank/credit/atm/AtmRequest.java}]{../../credit-service/src/main/java/by/bsuir/bank/credit/atm/AtmRequest.java}
\lstinputlisting[style=java, style=appendix, caption={credit-service/src/main/java/by/bsuir/bank/credit/atm/AtmResponse.java}]{../../credit-service/src/main/java/by/bsuir/bank/credit/atm/AtmResponse.java}
\lstinputlisting[style=java, style=appendix, caption={credit-service/src/main/java/by/bsuir/bank/credit/atm/DepositApi.java}]{../../credit-service/src/main/java/by/bsuir/bank/credit/atm/DepositApi.java}
\lstinputlisting[style=java, style=appendix, caption={credit-service/src/main/java/by/bsuir/bank/credit/web/AtmController.java}]{../../credit-service/src/main/java/by/bsuir/bank/credit/web/AtmController.java}
\lstinputlisting[style=java, style=appendix, caption={credit-service/src/test/java/by/bsuir/bank/credit/AtmProtocolIT.java}]{../../credit-service/src/test/java/by/bsuir/bank/credit/AtmProtocolIT.java}
\lstinputlisting[style=java, style=appendix, caption={credit-service/src/test/java/by/bsuir/bank/credit/CardNumbersTest.java}]{../../credit-service/src/test/java/by/bsuir/bank/credit/CardNumbersTest.java}

\subsection*{Запуск сервисов}

\lstinputlisting[style=shell, style=appendix, caption={scripts/run-all.sh}]{../../scripts/run-all.sh}
\lstinputlisting[style=shell, style=appendix, caption={scripts/stop-all.sh}]{../../scripts/stop-all.sh}


\end{document}
